% !TEX program = xelatex
% =====================================================================
%  Beamer template — adapted from weibo-hot/paper/pre/journal.club
%  Compile with XeLaTeX + Biber:  latexmk -xelatex main.tex
% =====================================================================
% \documentclass[draft]{beamer}
\documentclass[]{beamer}
\usetheme{Berlin}
\usepackage{hyperref}
\usepackage{graphicx}   % figures
\usepackage{multicol}   % multi-column layouts
\usepackage{booktabs}   % tables
\usepackage{subfigure}  % sub-figures
\usepackage{xcolor}
\usepackage{color}
\usepackage{csquotes}
\usepackage{xltxtra}
\usepackage{xelatexemoji}
\usepackage[style=apa,sortcites=true,sorting=nyt,backend=biber]{biblatex}
\addbibresource{refs.bib}
\usepackage{tikz}       % DAG graphs
\usetikzlibrary{shapes,decorations,arrows,calc,arrows.meta,fit,positioning}

\setbeamertemplate{caption}{\raggedright\insertcaption\par}
\setbeamertemplate{footline}{}
\tikzset{
  -Latex,auto,node distance =1 cm and 1 cm,semithick,
  state/.style ={ellipse, draw, minimum width = 0.7 cm},
  block/.style ={rectangle, draw, minimum width = 0.7 cm},
  point/.style = {circle, draw, inner sep=0.04cm,fill,node contents={}},
  bidirected/.style={Latex-Latex},
  el/.style = {inner sep=2pt, align=left, sloped}
}

% ---------- Colours ----------
\definecolor{green}{RGB}{62, 138, 46}
\definecolor{blue}{RGB}{3, 105, 170}
\definecolor{red}{RGB}{194, 54, 55}

\usecolortheme{beaver}

\definecolor{primary}{RGB}{0,24,49}
\definecolor{secondary}{RGB}{0,48,96}
\definecolor{tertiary}{RGB}{94,48,0}
\definecolor{quaternary}{RGB}{94,48,0}
\setbeamercolor{titlelike}{parent=palette secondary,fg=white}
\setbeamercolor{frametitle}{bg=primary!60!white}
\setbeamercolor{frametitle right}{bg=primary!60!white}
\setbeamercolor*{palette primary}{bg=primary, fg = white}
\setbeamercolor*{palette secondary}{bg=secondary, fg = white}
\setbeamercolor*{palette tertiary}{bg=tertiary, fg = white}
\setbeamercolor*{palette quaternary}{bg=quaternary, fg = white}

\definecolor{nus_blue}{rgb}{0,0.24,0.49}
\definecolor{nus_yellow}{rgb}{0.94, 0.48, 0.00}
\definecolor{nus_yellow_highlight}{rgb}{0.94, 0.48, 0.00}

\setbeamercolor{section in toc}{fg=black,bg=white}
\setbeamercolor{alerted text}{fg=nus_blue!80!gray}
\setbeamercolor{navigation symbols}{fg=gray!60!black, bg=gray!10!white}
\setbeamercolor{navigation symbols dimmed}{fg=gray!50!black}
\setbeamercolor*{palette primary}{fg=black!60!black,bg=gray!15!white}
\setbeamercolor*{palette secondary}{fg=black!70!black,bg=gray!15!white}
\setbeamercolor*{palette tertiary}{bg=black!80!black,fg=gray!10!white}
\setbeamercolor*{palette quaternary}{fg=black,bg=gray!5!white}
\setbeamercolor{enumerate item}{fg=black,bg=white}
\setbeamercolor{itemize item}{fg=black,bg=white}
\setbeamercolor{item}{fg=black,bg=white}
\setbeamercolor*{sidebar}{fg=black,bg=gray!15!white}

\setbeamercolor*{palette sidebar primary}{fg=black!10!black}
\setbeamercolor*{palette sidebar secondary}{fg=white}
\setbeamercolor*{palette sidebar tertiary}{fg=black!50!black}
\setbeamercolor*{palette sidebar quaternary}{fg=gray!10!white}

\setbeamercolor{button}{bg=white,fg=black}
\setbeamercolor{author in head/foot}{bg=black}
\setbeamercolor{title in head/foot}{bg=black!20}
\setbeamercolor{titlelike}{bg=black!100!white}
\setbeamercolor{frametitle}{bg=black!100!white}
\setbeamercolor{frametitle right}{bg=black!100!white}
\setbeamercolor{empty}{bg=black}

% ---------- Headline (section navigation) & frame title ----------
\setbeamertemplate{headline}{%
  \begin{beamercolorbox}[wd=\paperwidth,ht=3.5ex,dp=3.5ex]{section in head/foot}
    \insertnavigation{\paperwidth}%
  \end{beamercolorbox}%
}

\setbeamertemplate{frametitle}{%
  \nointerlineskip%
  \begin{beamercolorbox}[wd=\paperwidth,leftskip=.3cm,rightskip=.3cm plus1fil,vmode]{frametitle}
    \vskip0.5ex%
    \usebeamerfont*{frametitle}\strut\insertframetitle\strut\par%
    \vskip0.05ex%
  \end{beamercolorbox}%
}

\setbeamercolor*{separation line}{}
\setbeamercolor*{fine separation line}{}
\setbeamercolor{bibliography entry author}{fg=black}
\setbeamercolor{bibliography entry title}{fg=black}
\setbeamercolor{bibliography entry location}{fg=black}
\setbeamercolor{bibliography entry note}{fg=black}
\setbeamertemplate{bibliography item}{}

% ---------- CJK support ----------
\usepackage{xeCJK}

% ---------- Helper macros ----------
\providecommand{\tabfnm}[1]{\textsuperscript{#1}}  % apa7 marker used by \input tables
\newcommand{\hlt}[1]{\textcolor{red}{\textbf{#1}}}
\newcommand{\citegray}[1]{{\tiny\textcolor{gray}{\parencite{#1}}}}
\newcommand{\referencesize}{\tiny}
\renewcommand*{\bibfont}{\referencesize}

% Black divider slide. \sectionframe{Title} also starts a \section, so the
% headline navigation highlights correctly. Use \dividerframe{Title} for a
% divider without a new section (e.g. Acknowledgments, References).
\newcommand{\dividerframe}[1]{%
  {\setbeamercolor{background canvas}{bg=black}%
   \begin{frame}
     \begin{center}\Huge{\textcolor{white}{#1}}\end{center}
   \end{frame}}%
}
\newcommand{\sectionframe}[1]{\section{#1}\dividerframe{#1}}

% ---------- Title page ----------
\title[AI in Quantitative Social Research]{Using AI Responsibly in Quantitative Social Research}
\author[Rongxin Ouyang]{Rongxin Ouyang$^{*}$}
\institute{\textcolor{gray}{
    $^{*}$ PhD Candidate in Computational Communication, National University of Singapore; rongxin@u.nus.edu\\
Research Interest: Using causal and computational methods to evaluate societal and political impacts of digital information technologies}}
\date{\small \textcolor{gray}{November 5, 2026}}

\begin{document}

\frame{\titlepage}

% =====================================================================
\sectionframe{AI: A Bigger Picture}
% =====================================================================

\begin{frame}
  \frametitle{What Is AI?}
  \begin{columns}[T,onlytextwidth]
    \column{0.55\textwidth}
      \begin{itemize}
        \item ``The science and engineering of making intelligent machines'' \textcolor{gray}{\tiny (McCarthy)}
        \item Systems that \hlt{perceive} their environment and \hlt{act} to achieve goals
        \item A nested family of techniques, not a single technology
      \end{itemize}
    \column{0.45\textwidth}
      \begin{center}
        \begin{tikzpicture}[-,every node/.style={font=\tiny}]
          \draw[fill=gray!10] (0,0) ellipse (2.2 and 1.6);
          \draw[fill=gray!25] (0,-0.3) ellipse (1.7 and 1.2);
          \draw[fill=gray!40] (0,-0.6) ellipse (1.2 and 0.8);
          \draw[fill=red!30]  (0,-0.85) ellipse (0.7 and 0.45);
          \node at (0,1.3)   {Artificial Intelligence};
          \node at (0,0.6)   {Machine Learning};
          \node at (0,-0.1)  {Deep Learning};
          \node at (0,-0.85) {LLMs};
        \end{tikzpicture}
      \end{center}
  \end{columns}
\end{frame}

\begin{frame}
  \frametitle{A Brief History of AI}
  \footnotesize
  \begin{itemize}
    \item \textbf{1950s} Turing test (1950); ``AI'' coined at the Dartmouth workshop (1956)
    \item \textbf{1960s--80s} Symbolic AI and expert systems; two ``AI winters''
    \item \textbf{1990s--2000s} Statistical machine learning: learning rules from data
    \item \textbf{2012} Deep learning breakthrough in image recognition
    \item \textbf{2017} Transformer architecture
    \item \textbf{2018--2020} Pre-trained language models (BERT, GPT-2, GPT-3)
    \item \textbf{2022--} ChatGPT; \hlt{general-purpose LLMs} enter everyday research practice
  \end{itemize}
\end{frame}

% =====================================================================
\sectionframe{Pre-LLM AI for Quantitative Research}
% =====================================================================

\begin{frame}
  \frametitle{Pre-LLM: Machine Learning as a Measurement Tool}
  \begin{itemize}
    \item \textbf{Supervised learning}: hand-label a sample, train a classifier, scale up
    \item \textbf{Unsupervised learning}: topic models, clustering, word embeddings
    \item \textbf{Text, images, and networks as data}~\citegray{king2013how}
    \item \textbf{ML for causal inference}: heterogeneous effects, high-dimensional controls
  \end{itemize}
  ~\\
  \textcolor{gray}{Common thread: task-specific models, labeled data, and a researcher-built pipeline}
\end{frame}

% =====================================================================
\sectionframe{LLMs for Quantitative Research}
% =====================================================================

\begin{frame}
  \frametitle{Three Roles of LLMs in Research}
  \begin{center}
    \begin{tikzpicture}[-Latex]
      \node[block, minimum width=2.4cm, minimum height=0.9cm] (s) at (0,0)   {\small \textbf{Subject}};
      \node[block, minimum width=2.4cm, minimum height=0.9cm] (m) at (4,0)   {\small \textbf{Method}};
      \node[block, minimum width=2.4cm, minimum height=0.9cm] (o) at (8,0)   {\small \textbf{Object}};
      \node[font=\tiny, align=center] at (0,-1) {LLM \emph{stands in}\\for humans};
      \node[font=\tiny, align=center] (ml) at (4,-1) {LLM \emph{does}\\research work};
      \node[font=\tiny, align=center] at (8,-1) {LLM \emph{is studied}};
      \node[state, font=\tiny] (a) at (1.9,-2.6) {Assistant};
      \node[state, font=\tiny] (c) at (4,-2.6)   {Collaborator};
      \node[state, font=\tiny] (e) at (6.1,-2.6) {Evaluator};
      \path (ml) edge (a) (ml) edge (c) (ml) edge (e);
    \end{tikzpicture}
  \end{center}
\end{frame}

\begin{frame}
  \frametitle{LLM as Subject}
  \begin{itemize}
    \item \textbf{Silicon samples}: simulated survey respondents conditioned on personas
    \item \textbf{Simulated experiments}: pilot designs and power checks before fielding
    \item \textbf{Agent-based simulation}: populations of LLM agents interacting
  \end{itemize}
  ~\\
  \textcolor{gray}{Key question: when do simulated responses track real human behavior?}
\end{frame}

\begin{frame}
  \frametitle{LLM as Method: Assistant}
  \begin{itemize}
    \item Writing, debugging, and translating analysis code
    \item Searching and summarizing literature
    \item Editing and polishing prose
    \item Data cleaning and wrangling
  \end{itemize}
  ~\\
  \textcolor{gray}{The researcher specifies the task; the LLM speeds up execution}
\end{frame}

\begin{frame}
  \frametitle{LLM as Method: Collaborator}
  \begin{itemize}
    \item Brainstorming research questions and hypotheses
    \item Critiquing research designs and identification strategies
    \item Proposing measures, robustness checks, and alternative explanations
    \item Agentic workflows: multi-step analysis with researcher oversight
  \end{itemize}
  ~\\
  \textcolor{gray}{The LLM contributes ideas; the researcher remains accountable}
\end{frame}

\begin{frame}
  \frametitle{LLM as Method: Evaluator}
  \begin{itemize}
    \item \textbf{Annotation}: zero- and few-shot text classification at scale
    \item \textbf{LLM-as-judge}: scoring quality, stance, or similarity
    \item \textbf{Review}: feedback on manuscripts, code, and reproducibility
  \end{itemize}
  ~\\
  \textcolor{gray}{Must be validated against human-coded benchmarks}
\end{frame}

\begin{frame}
  \frametitle{LLM as Object}
  \begin{itemize}
    \item \textbf{Auditing models}: political slant, social bias, and value alignment
    \item \textbf{Effects on people}: persuasion, learning, labor, and well-being
    \item \textbf{Effects on information ecosystems}: AI-generated content and search
  \end{itemize}
\end{frame}

% =====================================================================
\sectionframe{Limitations}
% =====================================================================

\begin{frame}
  \frametitle{Limitations}
  \begin{enumerate}
    \item \textbf{Validity.} Outputs may not measure what we intend
    \item \textbf{Reproducibility.} Closed models, silent updates, and stochastic outputs
    \item \textbf{Bias.} Training data over-represent some groups and languages
    \item \textbf{Hallucination.} Fabricated facts and citations
    \item \textbf{Ethics.} Privacy, consent, and data sent to third parties
    \item \textbf{Over-reliance.} Deskilling and unexamined defaults
  \end{enumerate}
\end{frame}

% =====================================================================
\sectionframe{Looking Forward}
% =====================================================================

\begin{frame}
  \frametitle{Looking Forward: Responsible Use}
  \begin{itemize}
    \item \textbf{Disclose}: model, version, prompts, parameters, and date
    \item \textbf{Validate}: compare against human benchmarks before scaling
    \item \textbf{Reproduce}: pin versions; prefer open-weight models when possible
    \item \textbf{Stay accountable}: humans own the design, inference, and claims
  \end{itemize}
  ~\\
  \textcolor{gray}{From AI as a tool to AI as a research partner---with guardrails}
\end{frame}


{
  \setbeamercolor{background canvas}{bg=black}
  \begin{frame}
    \begin{center}
      \Huge{\textcolor{white}{Acknowledgments}}
    \end{center}
    \vspace{1cm}
    \normalsize{\textcolor{white}{We thank \ldots. All errors are our own.}}
  \end{frame}
}

\dividerframe{References}

\begin{frame}[allowframebreaks]
  \frametitle{References}
  \begingroup
  \referencesize\color{black}
  \printbibliography[notkeyword=chinese,heading=none]
  \endgroup
\end{frame}


\end{document}
